\section{Lecture 13: Why Deep Learning?}

\subsection{Motivation}

So far, we have studied:

\begin{itemize}
    \item linear and logistic regression,
    \item kernel methods and support vector machines,
    \item optimization and generalization principles.
\end{itemize}

\medskip

\noindent
These methods are powerful, but they rely heavily on \emph{manually designed features}.

\medskip

\noindent
\textbf{Key question.}  
Can we automatically learn representations from data?

\medskip

\noindent
\textbf{Guiding idea.}  
Deep learning combines representation learning with function approximation.

\subsection{Limitations of Shallow Models}

Classical models typically take the form:
\[
f(x) = w^\top \phi(x),
\]
where $\phi(x)$ is a fixed feature map.

\medskip

\noindent
\textbf{Limitations.}

\begin{itemize}
    \item \textbf{Feature engineering:}
    \begin{itemize}
        \item Requires domain expertise
        \item Difficult to scale
    \end{itemize}

    \item \textbf{Expressivity constraints:}
    \begin{itemize}
        \item May require very high-dimensional feature spaces
    \end{itemize}

    \item \textbf{Lack of hierarchy:}
    \begin{itemize}
        \item Cannot naturally capture compositional structure
    \end{itemize}
\end{itemize}

\medskip

\noindent
\textbf{Insight.}  
Shallow models separate representation from learning.

\subsection{Compositional Structure of Functions}

Many real-world functions exhibit hierarchical structure.

\medskip

\noindent
\textbf{Example.}

\begin{itemize}
    \item Images:
    \begin{itemize}
        \item edges $\rightarrow$ shapes $\rightarrow$ objects
    \end{itemize}

    \item Language:
    \begin{itemize}
        \item characters $\rightarrow$ words $\rightarrow$ sentences
    \end{itemize}
\end{itemize}

\medskip

\noindent
\textbf{Mathematical structure.}
\[
f(x) = f_L(f_{L-1}(\cdots f_1(x)\cdots)).
\]

\medskip

\noindent
\textbf{Key idea.}
\begin{itemize}
    \item Complex functions can be expressed as compositions of simpler functions
\end{itemize}

\medskip

\noindent
\textbf{Insight.}  
Hierarchical structure suggests using layered models.

\subsection{Motivation for Deep Learning}

Deep learning models take the form:
\[
f(x) = f_L \circ f_{L-1} \circ \cdots \circ f_1(x).
\]

\medskip

\noindent
\textbf{Advantages.}

\begin{itemize}
    \item \textbf{Representation learning:}
    \begin{itemize}
        \item Features are learned automatically
    \end{itemize}

    \item \textbf{Expressivity:}
    \begin{itemize}
        \item Deep models can represent complex functions efficiently
    \end{itemize}

    \item \textbf{Parameter sharing:}
    \begin{itemize}
        \item Reduces complexity in structured problems
    \end{itemize}
\end{itemize}

\medskip

\noindent
\textbf{Comparison with kernel methods.}

\begin{itemize}
    \item Kernels:
    \begin{itemize}
        \item Fixed feature space
    \end{itemize}

    \item Deep learning:
    \begin{itemize}
        \item Learned feature space
    \end{itemize}
\end{itemize}

\medskip

\noindent
\textbf{Insight.}  
Deep learning replaces manual feature design with learned representations.

\subsection{Transition to Learned Representations}

We reinterpret learning as:

\[
f(x) = w^\top \phi(x)
\quad \longrightarrow \quad
f(x) = w^\top \phi_\theta(x),
\]
where $\phi_\theta$ is learned from data.

\medskip

\noindent
\textbf{Key shift.}

\begin{itemize}
    \item Classical ML:
    \begin{itemize}
        \item Fixed $\phi(x)$
    \end{itemize}

    \item Deep learning:
    \begin{itemize}
        \item Learn $\phi_\theta(x)$ jointly with $w$
    \end{itemize}
\end{itemize}

\medskip

\noindent
\textbf{Optimization problem.}
\[
\min_{\theta, w} \; \sum_{i=1}^n \ell(w^\top \phi_\theta(x_i), y_i).
\]

\medskip

\noindent
\textbf{Challenge.}
\begin{itemize}
    \item Highly non-convex optimization
\end{itemize}

\medskip

\noindent
\textbf{Opportunity.}
\begin{itemize}
    \item Rich, task-specific representations
\end{itemize}

\medskip

\noindent
\textbf{Insight.}  
Learning representations and predictors jointly increases modeling power.

\subsection{A Unifying Perspective}

Deep learning builds on earlier ideas:

\begin{itemize}
    \item \textbf{Function approximation:} complex hypothesis spaces
    \item \textbf{Optimization:} gradient-based training
    \item \textbf{Generalization:} controlled through implicit and explicit regularization
    \item \textbf{Representation:} learned feature maps
\end{itemize}

\medskip

\noindent
\textbf{Core idea.}  
Deep learning integrates representation and prediction into a single optimization problem.

\medskip

\noindent
\textbf{Key insight.}  
The power of deep learning comes from composing simple transformations into complex structures.

\subsection{Outlook}

In this lecture, we examined:

\begin{itemize}
    \item limitations of shallow models,
    \item compositional structure of real-world functions,
    \item motivation for deep learning,
    \item transition to learned representations.
\end{itemize}

\medskip

\noindent
We now begin the study of neural networks.

\medskip

\noindent
In the next lecture, we introduce neural networks as compositional models and study their expressive power.

\medskip

\noindent
\textbf{Guiding principle.}  
Learning good representations is as important as learning the final predictor.